\section{Hardware Side-Channels and Pentimenti}
\label{sec:bg-sidechannels}

Hardware side-channels are unintended information channels that arise
from the physical implementation of computing systems
\cite{spreitzer2017systematic, hu2020overview}. The most well-studied of
these vulnerabilities are cache side-channel attacks which exploit CPU
microarchitectures allowing information to leak between users from
shared caches in multi-tenant environments~\cite{yarom2014flush,
	lipp2018meltdown, kocher2020spectre}. Generally, hardware side-channels
exploit physical phenomena that correlate with sensitive on-chip
activity~\cite{wei2020scoverview, nordholt2020hardware}. Classical
examples include power-analysis attacks that recover cryptographic keys
from power measurements or electromagnetic
measurements~\cite{kocher1999differential, mangard2007power,
	clavier2014practical, quisquater2001electromagnetic, agrawal2002side},,
and timing attacks that exploit data-dependent execution
latency~\cite{kocher1996timing, brumley2005remote}. Each class shares
the property that sensitive information is transmitted through a leakage
channel that the device manufacturer did not intend to expose, and each
has motivated a corresponding body of countermeasure research~\cite{
	barthe2018secure, ahmadi2023shield, zhang2024timing}.

The growth of cloud and edge deployments has substantially expanded the
attack surface for hardware side-channels. In multi-tenant cloud
infrastructure, users share physical hardware either concurrently
(\emph{spatial} co-tenancy) or sequentially over time (\emph{temporal}
co-tenancy)~\cite{ristenpart2009hey, zhang2014cloud, zhang2012crossvm}.
Spatial co-tenancy on shared FPGAs has been exploited extensively
through on-chip voltage-fluctuation sensors that recover sensitive
computation from co-located tenants~\cite{schellenberg2018remote,
	zhao2018fpga, gnad2017voltage, gnad2016analysis, schellenberg2018inside,
	ramesh2018fpga, giechaskiel2018leaky, giechaskiel2019leakier}. These
sensors form the basis of remote power-analysis attacks against
cryptographic accelerators and machine-learning workloads sharing the
fabric~\cite{schellenberg2018remote, gobulukoglu2021dac, kim2019make,
	mukhtar2018machine, bnn_power_side_channel}. Active denial-of-service
attacks have also been demonstrated through power-wasting circuits that
fault or reset co-tenant workloads~\cite{provelengios2020power,
	krautter2018fpgahammer, krautter2019active, krautter2019mitigating,
	nassar2021loopbreaker}. Edge platforms, with their more direct physical
accessibility, face analogous threats to embedded sensitive IP and user
data~\cite{tramer2016stealing, clements2022hardware, hu2016hardware,
	lu2020comprehensive}.

Temporal co-tenancy was recognized as a distinct threat more recently.
The defining characteristic of a temporal side-channel is that no
interaction between victim and adversary occurs during the victim's
execution: information leaks through residual physical state left behind
after the victim relinquishes the hardware. Existing demonstrations have
exploited thermal signatures that persisted across tenant
transitions~\cite{tian2019temporal}, residual charge in DRAM
cells~\cite{geist2020dram}, and BTI-induced wearout in FPGA routing
resources~\cite{drewes2024pentimento}. The last of these falls within
the attack class addressed in this dissertation.

Pentimenti are residual signatures left in the physical state of CMOS
transistors after they have carried sensitive
information~\cite{drewes2024pentimento}. The term is borrowed from the
art-historical practice of revealing earlier compositions beneath
finished paintings through X-ray or infrared imaging: just as a
painting's surface can hide traces of prior brushstrokes, a
semiconductor device's logical state can hide traces of prior internal
data. The mechanism underlying pentimenti is bias temperature
instability (BTI), a recoverable wearout phenomenon described in
Section~\ref{sec:bg-bti}. With sufficiently sensitive on-chip timing
instrumentation, the BTI-induced shifts in transistor switching time can
be measured and inverted to recover information about the bias history
of individual transistors, and the data those transistors carried. Prior
demonstrations of pentimenti recovery~\cite{drewes2024pentimento} have
established the attack only in the reconfigurable routing logic of
Xilinx FPGAs, where high-resolution phase-shift primitives provide the
timing precision required (Section~\ref{sec:bg-tdc}). Whether the attack
class and measurement techniques extend to other FPGA vendors, and more
broadly to other semiconductor device families whose physics involve
analogous defect dynamics, is the central question motivating this
dissertation.